% Opening of the old chapter c31, its \chapter command removed; generated
\label{chap:cx}

Chapter~\ref{chap:lo} turned the backward Cauchy flow into a positive linear inverse problem and solved it by the EM algorithm of Vardi and Lee. This chapter asks three questions about that computation. What does it cost, and what does log concave sampling in the sense of Applegate, Kannan and Polson say about it? Can EM be accelerated, by Nesterov momentum or otherwise, and what does the accelerated iteration compute? Does the representation of the Cauchy kernel as a scale mixture of normals help?

The short answers. The cost splits into three layers. Evaluating the level one measure needs $\zeta$ on the line $\Re s=1$, for which the best algorithms are polynomial in the height $T$ and not in $\log T$; this is the bottleneck, and the Thorin formulation does nothing to it. The inversion, given the data, is a convex minimisation whose Bayesian version is log concave sampling; the ill posedness appears as a precision requirement linear in the resolution. Verification is Turing's method and costs $T^{1+\varepsilon}$. EM is majorise minimise in the Kullback Leibler geometry; Nesterov momentum in the logarithmic domain keeps positivity, is an order of magnitude faster at $\theta'=1/2$ in the runs, and in every run reached the reverse information projection of the data onto the cone of Cauchy blurred positive measures, which is the Thorin measure at level $\theta'$ when it exists. The scale mixture of normals turns the Cauchy blur into a heat flow at a first passage time and identifies the Thorin measure with the subordinated heat trace of the zeros; it changes nothing in the complexity classification or the arithmetic.

The chapter then turns to the Thorin measures themselves: the levels $1$, $7/8$, $3/4$ and $1/2$, what positivity means, an integer valued mass deficit that decides positivity up to any finite height, a certificate in ball arithmetic up to height $1419.9$, the limits of the computation, and the global minimum of the Thorin density at the origin, unconditionally at level one (Proposition~\ref{bp:prop:v1min}) and, for every $\theta\in(\tfrac12,1]$, equivalent at all levels in $[\theta,1]$ to the zero free half plane $\Re s>\theta$ (Theorem~\ref{thm:dbeta:minall}). Its last part works open problems: an a posteriori certificate for the backward flow, a smoothed prime sum that evaluates $u_1$ to fixed accuracy in time polynomial in $\log v$, given that the zeros near that height lie on the critical line (verified below $3\times10^{12}$; at every height conditionally on the announced half plane $\Re s>7/8$), and the heuristic reason, from the explicit formula read backwards, why $\mathfrak U_{1/2}$ (equal to $2\widetilde U_{1/2}$ of Theorem~\ref{lv:thm:semigroup}) cannot be computed efficiently by deconvolution.

The chapter develops a source chapter on complexity and acceleration of the backward flow by Nick Polson, together with the computations done for this book. The numbers were recomputed for this book (Section~\ref{cx:sec:accel}, Theorem~\ref{cx:thm:certified}, the global minimum certificates and Proposition~\ref{cx:prop:slope}); where a recomputation differed from the source the recomputed value is printed, and numbers not recomputed are marked as the author's.

Notation. A level $\theta$ is the real part of $s=\theta+iv$, and $u_\theta(v)=\pi^{-1}\Re(\xi'/\xi)(\theta+iv)$; in the notation of Chapter~\ref{chap:tf} this is $\rho_{\theta-1/2}(v)$. The Cauchy density is $P_a(x)=(a/\pi)/(a^2+x^2)$, with $P_0=\delta_0$. A zero is $\rho=\sigma_\rho+i\gamma_\rho$. The symmetric Thorin measure $\mathfrak U_\theta$ of equation~\eqref{lo:eq:Utheta} has density $\sum_\rho P_{\theta-\sigma_\rho}(v-\gamma_\rho)$ with no factor $\tfrac12$, so $\mathfrak U_\theta=2\widetilde U_\theta$ with the $\widetilde U_\theta$ of Theorem~\ref{lv:thm:semigroup}. Certified means computed in ball arithmetic, with the sentence saying what the certificate covers. The announced zero free half plane $\Re s>7/8$ (OpenAI \cite{OAI26a,OAI26b}, formalised in Lean, and Liu \cite{Liu26a}; not yet refereed) is never used as a theorem: a result that uses it is stated as conditional on it. The results and their status are these.
\begin{enumerate}
\item The matrix $A$ of the discretised deconvolution, with cell densities as unknowns, has the properties (a) to (e) of Section~\ref{cx:sec:problem} (proved), including a kernel of dimension exactly $2(N-M)$ when the unknown window is larger. The level one data must be supplied to absolute precision $e^{-(1-\theta')\Omega}$ for the linear inverse (Proposition~\ref{cx:prop:precision}, proved). The Poisson posterior is log concave (Theorem~\ref{cx:thm:logconcave}); the sampling clause is a cited theorem, conditional on a bound for the ratio of radii not proved here.
\item EM has the rate $D(x^*\|x^{(0)})/k$ in the consistent case (Theorem~\ref{cx:thm:emrate}, proved, with the repaired step) and local linear convergence in the interior (Theorem~\ref{cx:thm:linear}, proved; its injectivity hypothesis fails in the runs). Under three growth conditions the EM limits converge vaguely to $\mathfrak U_{\theta'}$ when $\zeta$ has no zeros in $\Re s>\theta'$, and the residual stays bounded away from zero when it has one (Proposition~\ref{cx:prop:what}, proved).
\item Log domain Nesterov momentum with restart and SQUAREM are an order of magnitude faster than EM at $\theta'=1/2$ on a window; on the full range ($202$ zeros below height $400$) and on the first thousand zeros ($994$ of $995$ zeros below $1415$ found as peaks) they locate the zeros to about $10^{-2}$ (Section~\ref{cx:sec:accel}, computed; the rates $O(1/k)$ and $O(1/k^2)$ and the identification of the limit are observed, not proved).
\item The von Neumann and Ulam estimator for the backward flow has infinite variance whenever $h<0.5221\,\pi(1-\theta')$ (Proposition~\ref{cx:prop:vnvar}, proved); the positive reverse walk is the EM step in expectation. The Cauchy kernel is a first passage mixture of Gaussians, giving the subordinated heat trace of the zeros and an inverse gamma data augmentation sampler (Proposition~\ref{cx:prop:subord}, proved; the sampler comparison is computed).
\item The level $7/8$ Thorin measure is computed by three routes; its identification as a Thorin measure is conditional on the announced half plane. The positivity theorem (Theorem~\ref{cx:thm:positivity}), the mass deficit identity (Proposition~\ref{cx:prop:deficit}) and the level mass function (Corollary~\ref{cx:cor:levelmass}) are proved; the deficit tables, the $U_{3/4}$ and $U_{1/2}$ tables and Turing's check in Thorin bookkeeping are computed. Decidability by level is proved for finite heights (Section~\ref{cx:sec:limits}).
\item Certified positivity (Theorem~\ref{cx:thm:certified}, ball arithmetic at $96$ and $160$ bits): $u_{3/4}$ and $u_{7/8}$ are positive on $[0,1420]$, every certified lower bound exceeding $4\times10^{-5}$; the argument change of $\xi$ gives exactly $1000$ zeros below height $1419.9$, all on the critical line and simple, which is known (Platt and Trudgian \cite{PT21}); the positivity is what the certificate adds.
\item The global minimum of the Thorin density is at the origin: at level $1$ this is Proposition~\ref{bp:prop:v1min} (Theorem~\ref{cx:thm:globalmin}(i) gives an independent computed proof); at level $7/8$ it is conditional on the announced half plane and computed, the least certified bound on $[13.5,1420]$ being $0.0360$ at $\theta=1$ and $0.0272$ at $\theta=7/8$, both at $v=17.49$ (Theorem~\ref{cx:thm:globalmin}); the Riemann hypothesis form is Theorem~\ref{thm:dbeta:minall}(vi) (Theorem~\ref{cx:thm:rhmin}). The uniform law and the lower bound $W\ge w^*$ for the excess (Theorems~\ref{cx:thm:uniform} and~\ref{cx:thm:lbl}) are computed, and conditional on the announced half plane at level $7/8$.
\item The slope of the level one density has certified infima $0.1758$ to $0.1793$ (Proposition~\ref{cx:prop:slope}, computed); under the Riemann hypothesis the limiting bracket $X$ has its minimum in the first gap at $v^*=17.4979\pm0.0001$, with $0.1838985\le X(v^*)\le0.1838988$ (computed, conditional on RH); that this is its infimum over $13.5\le|v|\le1418.9$ is the author's computation, not rerun.
\item The a posteriori bound for window masses and centroids of the backward flow is proved (Proposition~\ref{cx:prop:apost}); the certificate of Theorem~\ref{cx:thm:apostcert} is the author's computation, not rerun, and is conditional on the absence of zeros to the right of $\theta'$ (the Riemann hypothesis at $\theta'=\tfrac12$); it certifies Gaussian weighted masses and centroids, and single zeros only through Corollary~\ref{cx:cor:zeroloc} (proved).
\item The smoothed prime sum for $u_1$ has an exact identity (proved) and an error bound proved under the local hypothesis on the zeros, including the pole term the source omitted (Theorem~\ref{cx:thm:fastu1}); conditionally on the announced half plane it holds at every height with $X$ of order $(\log v/\varepsilon)^8$; the nonuniform FFT experiment is the author's, not rerun.
\item The explicit formula read backwards and the efficiency of the backward flow (Section~\ref{cx:sec:halfeff}), and the reading of the Hilbert and P\'olya conjecture as a resolvent trace (Section~\ref{cx:sec:vonneumann}), are heuristic.
\end{enumerate}
